% ============================================== %

%%%%%%%%%%%%%%%%%%%%%%%%%%%%%%%%%%%%%%%%%%%%%%%%%%
%%%			  07 Trabalho Prático 07		   %%%
%%%%%%%%%%%%%%%%%%%%%%%%%%%%%%%%%%%%%%%%%%%%%%%%%%

% ============================================== %

\section{Trabalho Prático 07}

\begin{enumerate}
    \item Calcular as energias de deformação específicas para as medidas de deformação
    Logarítmica e de Green.
    \item Encontrar os valores das constantes para leis do tipo: $\sigma^* = E\cdot
    \varepsilon^*$ (leis Hookeanas).
\end{enumerate}

% ================ Resolução TP-07 ================
\begin{textocaixa}
\hspace{10mm}

A energia de deformação específica $u^*$ se relaciona com o estiramento $\lambda$ a partir
da \autoref{eq:TP7-energia}.

\begin{equation}
    u^* = \int_{\varepsilon^*}\sigma^*\,\mathrm{d}\varepsilon^*
        = \int_\lambda \left(\sigma^*\:\frac{\partial \varepsilon^*}
                                            {\partial \lambda}\right)\,\mathrm{d}\lambda
    \label{eq:TP7-energia}
\end{equation}

Em consequência, considerando lei constitutiva Hookeana para os pares de tensão-deformação,
a expressão para o cálculo da energia de deformação específica se transforma em:

\begin{equation}
    u^* = E\:\int_\lambda \left(\varepsilon^*\:
                                \frac{\partial \varepsilon^*}
                                     {\partial \lambda}\right) \mathrm{d}\lambda
    \label{eq:TP7-energia_estiramento}
\end{equation}

Ou, então, numa relação direta com a deformação $\varepsilon^*$, dada na
\autoref{eq:TP7-energia_deform}, que será utilizada nesta resolução.

\begin{equation}
    u^* = E\:\int_{\varepsilon^*} \left(E\cdot\varepsilon^*\right)
                                      \,\varepsilon^* \mathrm{d}\varepsilon^*
    \longrightarrow \boxed{u^* = \frac{E}{2}\left[\varepsilon^*\left(\lambda\right)\right]^2}
    \label{eq:TP7-energia_deform}
\end{equation}

\begin{enumerate}
    \item \underline{Energia de deformação específica}

A medida de deformação Logarítmica $\varepsilon^* = \varepsilon_L$ é expressa por:

\begin{equation}
    \varepsilon_L = \ln \lambda
    \label{eq:TP7-eL}
\end{equation}

Implicando que a energia de deformação para a medida de deformação Logarítmica é dada por:

\begin{equation}
    \boxed{u_L = \frac{E}{2}\left(\ln\lambda\right)^2}
    \label{eq:TP7-uL_general}
\end{equation}

Já a medida de deformação de Green se expressa na \autoref{eq:TP7-eG}, permitindo desenvolver
a expressão para a energia de deformação específica para esta medida de deformação ($u_G$)
na \autoref{eq:TP7-uG}.

\begin{equation}
    \varepsilon_G = \frac{1}{2}\left(\lambda^2 - 1\right)
    \label{eq:TP7-eG}
\end{equation}

\begin{equation}
    \Longrightarrow u_G = \frac{E}{2}\left[\frac{1}{2}\left(\lambda^2 - 1\right)\right]^2
    \longrightarrow
    \boxed{u_G = \frac{E\,\lambda^2}{8}\left(\lambda^2 - 2\right) + \frac{E}{8}}
    \label{eq:TP7-uG}
\end{equation}

\item \underline{Constantes para leis do tipo Hookeanas}

Os conjugados de tensão de cada medida da deformação clássica se relacionam com a tensão de
Biot $\sigma_B$, dada na \autoref{eq:TP7-sigmaB}, a partir de expressões envolvendo o
estiramento.

\begin{equation}
    \sigma_B = E\left(\lambda - 1\right)
    \label{eq:TP7-sigmaB}
\end{equation}

Para as medidas de deformação Logarítmica e de Green, essas relações são mostradas na
\autoref{eq:TP7-sigmaLGB}.

\begin{bigEquation}
    \left\{\begin{array}{lclcl}
        \sigma_L &=& \lambda\,\sigma_B &=& 
                     E\cdot \left(\lambda^2 - \lambda\right) \\ [10pt]
        \sigma_G &=& \frac{1}{\lambda}\,\sigma_B &=& 
                     E\cdot \left(1 -\frac{1}{\lambda}\right) \\
    \end{array}\right.
    \label{eq:TP7-sigmaLGB}
\end{bigEquation}

Além disso, considerando a constante de proporcionalidade $E_L$ para a medida de deformação
Logarítmica e $E_G$ para a medida de deformação de Green, as leis constitutivas equivalente
são expressas como:

\begin{bigEquation}
    \left\{\begin{array}{lclcl}
        \sigma_L &=& E_L\cdot\varepsilon_L &=& 
                     E_L\cdot \ln\lambda \\ [10pt]
        \sigma_G &=& E_G\cdot\varepsilon_G &=& 
                     E_G\cdot \frac{1}{2}\left(\lambda^2 - 1\right) \\
    \end{array}\right.
    \label{eq:TP7-sigmaLG_const}
\end{bigEquation}

Portanto, comparando-se a \autoref{eq:TP7-sigmaLG_const} com a \autoref{eq:TP7-sigmaLGB},
encontram-se relações entre as constantes de proporcionalidade e o módulo de Young $E$
referente à medida de deformação de Biot -- a qual está ligada a ensaios de caracterização
que medem a deformação com base nas dimensões iniciais. Para a medida de deformação
Logarítmica, tem-se, pois, que:

\begin{equation}
    \frac{E_L}{E} = \frac{\lambda^2 - \lambda}{\ln\lambda}
    \label{eq:TP7-EL_E}
\end{equation}

Enquanto, para a medida de deformação de Green, a razão entre $E_G$ e $E$ se torna:

\begin{equation}
    \frac{E_G}{E} = \frac{2}{\lambda^2 + \lambda}
    \label{eq:TP7-EG_E}
\end{equation}

Retomando esse resultado na \autoref{eq:TP7-sigmaLGB}, as leis constitutivas do tipo
Hookeanas para as medidas de deformação Logarítmica e de Green são formuladas na
\autoref{eq:TP7-sigmaL_hook} e \autoref{eq:TP7-sigmaG_hook}, respectivamente.

\begin{equation}
    \boxed{\sigma_L = \frac{\lambda^2 - \lambda}{\ln\lambda}\,E\cdot \varepsilon_L}
    \label{eq:TP7-sigmaL_hook}
\end{equation}

\begin{equation}
    \boxed{\sigma_G = \frac{2E}{\lambda^2 + \lambda}\cdot \varepsilon_G}
    \label{eq:TP7-sigmaG_hook}
\end{equation}

\end{enumerate}



\end{textocaixa}